\section{Exact minimax values through three switches}\label{sec:small-minimax}

The analytic reach bounds and the four-block certificate theorem determine
both the one-sided minimax and the full minimax. A common formulation also
preserves the stronger conclusion available under equal terminal masses.

\begin{theorem}[One-sided minimax through four blocks]\label{thm:small-one-sided}
For $1\le k\le4$ and $n\ge k+1$,
\begin{equation}\label{eq:small-one-sided}
 G^-_{n,k}(T)=\frac{T}{n(r_n^k-1)}.
\end{equation}
The constant uniform input attains this value. For $k=4$ the upper bound
uses the computer-assisted Theorem~\ref{thm:reach-four}; for $k\le3$ its
proof is analytic.
\end{theorem}
\begin{proof}
Use the single-, two-, three-, or four-block reach result, respectively,
at $E=T/[n(r_n^k-1)]$. It reaches $T$ and proves the upper bound with
distinct modes. For the uniform input $A_i(t)=t/n$, let an arbitrary
schedule, allowing repetitions, have one-sided error $E$ and block
endpoints $t_j$. At the end of block $j$, the active mode's occupation
is at least $t_j-t_{j-1}$, so
\[
 t_j-t_{j-1}-t_j/n\le E,
 \qquad t_j\le r_n(t_{j-1}+E).
\]
Iteration yields $T\le nE(r_n^k-1)$. Shorter schedules can be padded
with zero-length blocks, or the recurrence can be iterated further using
its monotonicity. This proves the matching lower bound even when modes
repeat.
\end{proof}

\begin{theorem}[Full minimax through three switches]\label{thm:small-full}
For $0\le s\le3$ and $n\ge s+2$,
\begin{equation}\label{eq:small-full}
 F_{n,s}(T)=T\max\left\{\frac1{s+2},
                  \frac1{n(r_n^{s+1}-1)}\right\}.
\end{equation}
In particular,
\begin{align}
 F_{n,2}(T)&=T\max\left\{\frac14,
           \frac{(n-1)^3}{n(3n^2-3n+1)}\right\}, && n\ge4,
           \label{eq:two-switch-exact}\\
 F_{n,3}(T)&=T\max\left\{\frac15,
           \frac{(n-1)^4}{n(4n^3-6n^2+4n-1)}\right\}, && n\ge5.
           \label{eq:three-switch-exact}
\end{align}
The three-switch upper bound is computer-assisted; the other cases are
analytic.
\end{theorem}
\begin{proof}
Combine Theorems~\ref{thm:one-sided-reduction} and
\ref{thm:small-one-sided}. The two matching input types are the $s+2$
successive distinct pure modes of equal duration and the constant uniform
input. Theorem~\ref{thm:uniform} gives the latter's full error, whose
additional $T/n$ term is bounded by $T/(s+2)$ in this range.
\end{proof}

\begin{corollary}[Equal terminal masses]\label{cor:equal-masses}
Let $1\le k\le4$ and $n\ge k+1$. Over inputs with $m_i=T/n$ for every
mode, the exact worst-case full error with at most $k$ blocks is
\begin{equation}\label{eq:equal-masses}
 \max_{\substack{A\in\Acal_n(T)\\A_i(T)=T/n\ (i\in[n])}}
       \OPT_{k-1}(A)
 =T\max\left\{\frac1n,\frac1{n(r_n^k-1)}\right\}.
\end{equation}
The constant uniform input is an extremizer in this restricted class.
\end{corollary}
\begin{proof}
The reach theorem controls $W-A$ at the second threshold in
\eqref{eq:equal-masses}. Independently, every positive discrepancy is
bounded by $A_i(t)\le m_i=T/n$. Combining these two bounds proves the
upper bound. The uniform input belongs to the class and attains the
displayed value by Theorem~\ref{thm:uniform}.
\end{proof}

The equal-mass conclusion is stronger than merely applying the unrestricted
minimax bound to these inputs. For two switches it gives $T/4$ at $n=4$
and the uniform recurrence value for all $n\ge5$. Thus at $n=5,6,7$ the
restricted worst case is strictly smaller than the unrestricted value
$T/4$. For three switches it gives $T/n$ at $n=5,6$ and the uniform
recurrence value for $n\ge7$; it is strictly smaller than $T/5$ for
$n=6,\ldots,11$. The common corollary derives the equal-mass result for
four blocks as well as the earlier analytic three-block result without
needing a separate construction for each mass-constrained class.

\subsection{Plateau transitions and mode-count asymptotics}

For two switches, the uniform coefficient in~\eqref{eq:two-switch-exact}
exceeds $1/4$ exactly when
\[
 p_2(n)=n^3-9n^2+11n-4>0.
\]
Direct substitution gives negative values at $n=4,5,6,7$. At $n=8+h$,
\[
 p_2(8+h)=h^3+15h^2+59h+20>0\qquad(h\ge0).
\]
Thus $F_{n,2}(T)=T/4$ for $4\le n\le7$, and uniform controls are
worst-case for all $n\ge8$.

For three switches, the corresponding sign polynomial is
\[
 p_3(n)=n^4-14n^3+26n^2-19n+5.
\]
It is negative at $n=5,\ldots,11$, whereas
\[
 p_3(12+h)=h^4+34h^3+386h^2+1469h+65>0\qquad(h\ge0).
\]
Hence $F_{n,3}(T)=T/5$ for $5\le n\le11$, and uniform controls are
worst-case for $n\ge12$. Expanding the exact rational expressions gives
\begin{align}
 F_{n,2}(T)&=\frac T3-\frac{2T}{3n}+\frac{2T}{9n^2}
                      +O(Tn^{-3}),\label{eq:two-switch-expansion}\\
 F_{n,3}(T)&=\frac T4-\frac{5T}{8n}+\frac{5T}{16n^2}
                      +O(Tn^{-3}).\label{eq:three-switch-expansion}
\end{align}
These are expansions of the full minimax, since the uniform term dominates
for all sufficiently large $n$, rather than lower-bound expansions alone.

\section{Structural limits of the constructions}\label{sec:structural-limits}

The hypotheses and choices in the preceding arguments matter. We give
exact counterexamples to three tempting extensions. In the first example,
the failure extends beyond distinct-mode schedules: allowing a repeated
mode still does not restore uniform extremality at the boundary $k=n$.
This strengthens what a failure of a distinct-mode construction alone
would imply.

\subsection{Three modes and three blocks: uniform input need not be worst-case}

\begin{proposition}[Failure at $k=n=3$]\label{prop:three-mode-failure}
There is an input $A\in\Acal_3(57/8)$ with
$\OPT^-_3(A)>1$, although the constant uniform input on that horizon
has one-sided optimum exactly one with three blocks. Moreover, every
three-distinct-mode reach composition at threshold one stops at
$971/146<57/8$ for this input. In particular,
$G^-_{3,3}(57/8)>1$.
\end{proposition}
\begin{proof}
For this example label the modes $0,1,2$. Define their cumulative
allocations by linear interpolation through Table~\ref{tab:n3-knots},
dividing every entry by 146. After its last knot $t_*=971/146$, extend
with rates $(1/3,1/3,1/3)$ to $L=57/8$. Every allocation increment is
nonnegative and the three increments sum to its time increment. Each
individual slope is at most $3/4$, so this is an admissible control and
every $H_i(t)=t-A_i(t)$ is strictly increasing.

\begin{table}[tb]
\centering
\begin{tabular}{rrrr}
\toprule
$146t$ & $146A_0(t)$ & $146A_1(t)$ & $146A_2(t)$\\
\midrule
0&0&0&0\\
146&98&48&0\\
194&110&48&36\\
256&110&78&68\\
408&224&116&68\\
516&224&178&114\\
580&240&178&162\\
971&417&309&245\\
\bottomrule
\end{tabular}
\caption{Rational cumulative-allocation knots for the three-mode example.
Every entry, including time, is divided by 146.}
\label{tab:n3-knots}
\end{table}

At threshold $E=1$ the first reaches are
\[
 (R_0,R_1,R_2)=(128/73,97/73,1).
\]
Both orders of the other two modes after excluding $i$ have the same
pair reach, respectively
\[
 (M_0,M_1,M_2)=(204/73,258/73,290/73).
\]
These identities follow by substituting the relevant knots into
$H_i(R_i)=1$ and $H_k(M_i)=R_j+1$ for distinct $i,j,k$.
At the last knot, $H_i(t_*)=M_i+1$ for each $i$. Strict increase proves
that all six three-distinct-mode compositions end exactly at $t_*$.

It remains to rule out repeated modes; failure of these six words alone
would not establish the proposition. The terminal masses are
\begin{equation}\label{eq:n3-terminal}
 (m_0,m_1,m_2)=\frac1{1752}(5281,3985,3217),
 \qquad m_0>2,\quad m_1>2.
\end{equation}
For any one-sided error-one schedule with active-mode set $J$, its
terminal inequalities give
\[
 L=\sum_{i\in J}W_i(L)\le\sum_{i\in J}(m_i+1),
 \qquad \sum_{i\notin J}m_i\le |J|.
\]
If at most two modes are used, omitting mode 0 or 1 is impossible by
\eqref{eq:n3-terminal}. The only possible active set is therefore
$\{0,1\}$. Every schedule on this set with at most three blocks can be
written as $p,q,p$, with $(p,q)=(0,1)$ or $(1,0)$, allowing zero-length
blocks to include shorter schedules.

Let its first two endpoints be $u\le v$. The first two negative-error
constraints give $u\le R_p$ and $v\le\Phi_q(u)\le\Phi_q(R_p)=M_2$.
Moreover,
\begin{equation}\label{eq:middle-service}
 v-u\le A_q(v)+1\le A_q(M_2)+1=M_2-R_p.
\end{equation}
The final occupation of $p$ is $L-(v-u)$; its terminal constraint requires
$v-u\ge L-m_p-1$. For $p=0$ the upper bound in
\eqref{eq:middle-service} is $162/73$, whereas the required lower bound
is $2725/876$, exceeding it by $781/876$. For $p=1$ the corresponding
bounds are $193/73$ and $3373/876$, with gap $1057/876$. Both are
impossible. This rules out every schedule with at most three blocks at
error one. Attainment of the instance minimum
(Proposition~\ref{prop:compactness}) then gives the strict inequality
$\OPT^-_3(A)>1$.

Finally, the uniform input has the distinct-mode schedule with endpoints
$3/2,15/4,57/8$; each active mode has negative discrepancy one at the end
of its block. Conversely, the block-end recurrence
$t_j\le(3/2)(t_{j-1}+E)$, valid even with repeated modes, gives
$57/8\le3E((3/2)^3-1)=(57/8)E$. Thus its one-sided optimum is exactly
one, as asserted.
\end{proof}

The supplement checks the knot slopes, all six distinct reach compositions,
terminal masses, and both middle-service contradictions using rational
arithmetic. This counterexample does not concern the spare-mode range
$k<n$ of Theorems~\ref{thm:one-sided-reduction} and
\ref{thm:small-one-sided}, nor does it have equal terminal masses.

\subsection{A specified heavy mode need not admit an adjacent pair}

Theorem~\ref{thm:heavy} chooses a repeat; it does not promise to repeat
any prescribed heavy mode. On four unit intervals, consider the interval
rates
\[
 \begin{array}{c|cccc|c}
 \text{mode}&1&2&3&4&\text{total}\\\hline
 0&7/15&0&0&0&7/15\\
 1&2/15&1/2&0&4/5&43/30\\
 2&2/5&1/2&1&1/5&21/10
 \end{array}
\]
Mode 1 is heavy at threshold one. If an error-one unit-slot word selected
it in slots 1--2 or 2--3, its occupation at the pair's end would be at
least two, but its allocation is only $19/30$. This gives negative error
at least $41/30>1$. If its pair occupied slots 3--4, mode 2 would need
at least two selections in slots 1--2, because its terminal mass is
$21/10>2$. Yet two such selections give error $2-9/10=11/10>1$ at
prefix 2. No location of the prescribed pair works. The words $(0,2,2,1)$
and $(0,1,2,2)$ both have error at most one, consistently with the
existential theorem. This statement concerns rounding on the fixed unit
slots; it does not impose a restriction on arbitrary continuous block
lengths.

\subsection{The pair position cannot be fixed by the first unit-mass prefix}

Even if a largest-mass mode is chosen, the following rule can fail: force
its two adjacent slots to end at the first integer prefix where its
allocation reaches one. On five unit intervals take
\[
 \begin{array}{c|ccccc|c}
 \text{mode}&1&2&3&4&5&\text{total}\\\hline
 0&3/5&0&1/10&2/5&1&21/10\\
 1&3/10&3/5&11/20&3/5&0&41/20\\
 2&1/10&2/5&7/20&0&0&17/20
 \end{array}
\]
Mode 0 has strictly largest mass and first reaches allocation at least
one at prefix 4. Forcing it in slots 3--4 leaves mode 1 with no service
there. At prefix 4, mode 1 already has allocation $41/20>2$, so an
error-one word must have selected it in both slots 1--2. Its allocation
at prefix 2 is only $9/10$, giving the impossible negative error
$11/10$. In contrast, $(0,1,1,0,0)$ has error at most one and repeats
the largest mode in slots 4--5. Thus both the selected mode and the pair
position are meaningful choices in the prefix-reordering proof. These
examples do not disprove the separate conjecture that some error-one
word can always repeat a largest heavy mode adjacently; that strengthening
is not used by any theorem here.
